\section{Conclusion}
A piecewise Johnson--Cook constitutive framework with a phonon-drag-motivated high-rate increment has been reconstructed from the model-development logic established in the present study. The final form preserves the conventional low-rate JC branch, introduces a fixed effective transition strain rate tied to the material scales $c_{T0}$ and $b$ through the bridge factor $\Lambda_{\mathrm{ph}}$, and adds a logarithmic drag increment that is continuous at the transition.

Several modeling decisions were essential. First, the characteristic sound-speed scale was kept constant, while the temperature dependence was retained only in the shear-modulus factor $G(T)$. Second, the critical strain rate was treated as a material-dependent constitutive threshold rather than a direct function of impact velocity. Third, the high-rate branch was written as the continuing low-rate JC baseline plus a drag increment, rather than as a drag-only stress. This preserved continuity and avoided the nonphysical stress drop that appears when the drag term alone is used above the transition.

The resulting closure is not a full microscopic dislocation model, but it provides a compact and physically interpretable constitutive form for high-rate cold-spray impact simulations. It separates the roles of hardening, low-rate JC rate sensitivity, thermal softening, transition-rate bridging, and high-rate drag amplification in a way that can be implemented directly in the numerical solver and calibrated against impact data.
